\chapter{最適割当問題}
\label{ch:sem-assignment-chapter}

\begin{remark}{本章の設定}{sem-ch2-setup}
  $(\X, d_\X)$, $(\Y, d_\Y)$ は\textbf{完備可分な距離空間}（Polish 空間，
  定義~\ref{def:sem-polish}）とする．
  ボレル $\sigma$-代数 $\Bb(\X)$ により $\X$ を可測空間 $(\X, \Bb(\X))$ として扱う
  （$\Y$ についても同様）．$\X$ 上の確率測度全体を $\Mm_+^1(\X)$ と記す．

  Polish 空間の3要素を確認する：
  \begin{itemize}
    \item \textbf{距離空間}．距離関数
      $d_\X \colon \X \times \X \to [0, \infty)$ が距離の公理
      （定義~\ref{def:sem-metric-space}）を満たし，$d_\X$ から以下が順次定まる：
      \begin{align*}
        B_r(x)
          &\defeq \{y \in \X \mid d_\X(x, y) < r\}
          && \text{（開球，定義~\ref{def:sem-open-ball}）} \\
        \mathcal{O}_{d_\X}
          &\defeq \{U \subset \X \mid \forall x \in U,\,
            \exists r > 0,\, B_r(x) \subset U\}
          && \text{（開集合の全体）} \\
        \Bb(\X)
          &\defeq \sigma(\mathcal{O}_{d_\X})
          && \text{（ボレル } \sigma\text{-代数，定義~\ref{def:sem-borel}）}
      \end{align*}
      ここで $\sigma(\mathcal{O}_{d_\X})$ は，$\mathcal{O}_{d_\X}$ を含む
      $\X$ 上のすべての $\sigma$-代数の共通部分
      \[
        \sigma(\mathcal{O}_{d_\X})
        \defeq \bigcap
        \bigl\{
          \mathcal{F} \subset 2^{\X}
          \;\big|\;
          \mathcal{F} \text{ は } \sigma\text{-代数, }
          \mathcal{O}_{d_\X} \subset \mathcal{F}
        \bigr\}
      \]
      として定まる．これにより可測空間 $(\X, \Bb(\X))$
      （定義~\ref{def:sem-measurable-space}）が得られる．
    \item \textbf{完備性}．
      列 $(x_n)_{n \geq 1} \subset \X$ が\textbf{Cauchy 列}であるとは，
      任意の $\varepsilon > 0$ に対してある $N \in \N$ が存在して
      $n, m \geq N \Longrightarrow d_\X(x_n, x_m) < \varepsilon$ が成り立つことをいう．
      $(\X, d_\X)$ が完備とは，任意の Cauchy 列 $(x_n) \subset \X$ が
      $\X$ の点に収束する（定義~\ref{def:sem-convergence}）ことをいう．
      すなわち
      \[
        \forall (x_n) \subset \X \text{ Cauchy 列},\;
        \exists x \in \X,\; x_n \to x.
      \]
      収束の定義に展開すれば
      \[
        \forall (x_n) \subset \X \text{ Cauchy 列},\;
        \exists x \in \X,\;
        \forall \varepsilon > 0,\;
        \exists N \in \N,\;
        n \geq N \Longrightarrow d_\X(x_n, x) < \varepsilon.
      \]
    \item \textbf{可分性}．
      $D \subset \X$ が\textbf{稠密}であるとは，任意の空でない開集合
      $U \subset \X$ に対し $D \cap U \neq \emptyset$ となることをいう
      （定義~\ref{def:sem-dense}）．
      距離空間ではこれは「任意の $x \in \X$ と任意の $\varepsilon > 0$ に対し
      $\exists q \in D,\; d_\X(x, q) < \varepsilon$」と同値．
      $(\X, d_\X)$ が可分とは可算な稠密部分集合 $D \subset \X$ が
      存在することをいう（定義~\ref{def:sem-separable}）．すなわち
      \[
        \exists D \subset \X \text{ 可算},\;
        \forall U \subset \X \text{ 空でない開集合},\;
        D \cap U \neq \emptyset.
      \]
  \end{itemize}

  離散版では確率測度の表示に現れる点集合
  $S \defeq \{x_1, \ldots, x_n\} \subset \X$ は有限集合となる．
  $S$ にどの距離 $d$ を入れた距離空間 $(S, d)$ も
  自動で Polish 空間となる（$\X$ の可測構造は $S$ に制限することで継承される）：
  \begin{itemize}
    \item[] \textbf{完備性．} $n = 1$ なら $S = \{x_1\}$ ゆえ
      任意の列は $y_k = x_1$ ($\forall k$) なる定値列であり，
      $d(y_k, x_1) = 0$ より $x_1$ に収束する．
      以下 $n \geq 2$ とする．相異なる 2 点間の距離の最小値を
      $\delta \defeq \min_{i \neq j} d(x_i, x_j)$ とおく．
      これは有限個の正数の最小値なので $\delta > 0$．
      任意の Cauchy 列 $(y_k)_{k \geq 1} \subset S$ に対し，
      Cauchy 条件で $\varepsilon \defeq \delta / 2$ を取ると，
      ある $N \in \N$ が存在して
      $k, l \geq N \Longrightarrow d(y_k, y_l) < \delta / 2 < \delta$．
      一方 $S$ の相異なる 2 点の距離は $\delta$ 以上だから，
      $d(y_k, y_l) < \delta$ は $y_k = y_l$ を強制する．
      ゆえに列 $(y_k)$ は $N$ 以降すべて同一の点となり，その点に収束する．
    \item[] \textbf{可分性．} $S$ 自身が有限（特に可算）集合であり，
      自身の中で稠密だから可分．
  \end{itemize}
  したがって本章の Polish 強化は連続側で必要となる一方，
  離散版には追加制約を課さない．
\end{remark}

\begin{definition}{最適割当問題}{sem-assignment}
  $n \in \N$，置換 $\sigma \in \Perm(n)$，
  および行列 $\mathbf{C} \in \R_{\geq 0}^{n \times n}$ に対して，
  \[
    \min_{\sigma \in \Perm(n)} \frac{1}{n} \sum_{i=1}^{n} C_{i, \sigma(i)}
  \]
  を達成する $\sigma$ を求める問題を\textbf{最適割当問題}という．
\end{definition}

添字 $i$ から $j$ への\textbf{輸送コスト}を $C_{i,j}$ と解釈すれば，
$\sigma(i) = j$ は $i$ を $j$ に割り当てることを意味し，
目的関数 $\frac{1}{n}\sum_i C_{i, \sigma(i)}$ はその割当の平均輸送コストを表す．

\begin{example}{最適割当}{sem-assignment-cost-example}
  $n = 2$ とし，置換 $\sigma \in \Perm(2)$ と行列
  \[
    \mathbf{C} =
    \begin{pmatrix}
      1 & 4 \\
      2 & 1
    \end{pmatrix} \in \R_{\geq 0}^{2 \times 2}
  \]
  を考える．$\Perm(2) = \{(1,2),\, (2,1)\}$ の $2! = 2$ 通りの割当を以下に図示する．

  \begin{center}
  \begin{tikzpicture}[
    wk/.style={circle, draw=blue!60, fill=blue!10,
               minimum size=18pt, inner sep=1pt, font=\small},
    tk/.style={rectangle, rounded corners=1.5pt,
               draw=red!60, fill=red!10,
               minimum size=18pt, inner sep=1pt, font=\small},
    optarr/.style={-{Stealth[length=5pt]}, thick, green!50!black},
    arr/.style={-{Stealth[length=5pt]}, semithick, black!60, dashed},
    costlbl/.style={fill=white, inner sep=1pt, font=\footnotesize},
  ]
    \node[wk] (w1) at (0, 1.2) {1};
    \node[wk] (w2) at (3, 1.2) {2};
    \node[tk] (t1) at (1, -1.2) {1};
    \node[tk] (t2) at (4, -1.2) {2};

    % σ=(1,2): 1→1, 2→2（最適）
    \draw[optarr] (w1) -- node[costlbl, midway] {(a)\ $1$} (t1);
    \draw[optarr] (w2) -- node[costlbl, midway] {(a)\ $1$} (t2);
    % σ=(2,1): 1→2, 2→1
    \draw[arr] (w1) -- node[costlbl, pos=0.25] {(b)\ $4$} (t2);
    \draw[arr] (w2) -- node[costlbl, pos=0.25] {(b)\ $2$} (t1);
  \end{tikzpicture}
  \end{center}

  各割当のコストをまとめると，
  \begin{center}
  \begin{tabular}{ccc}
     & $\sigma$ & コスト $\frac{1}{n}\sum_i C_{i,\sigma(i)}$ \\\hline
    (a) & $(1,2)$ & $\frac{1}{2}(1 + 1) = \mathbf{1}$ \\
    (b) & $(2,1)$ & $\frac{1}{2}(4 + 2) = 3$ \\
  \end{tabular}
  \end{center}
  よって $\sigma = (1, 2)$ が最小コスト $1$ を達成し，最適である．
\end{example}

\begin{claim}{最適解の存在}{sem-assignment-existence}
  任意の $n \in \N$ と任意の $\mathbf{C} \in \R_{\geq 0}^{n \times n}$ に対して，
  最適割当問題の最小値は達成される．
\end{claim}

\begin{proof}
  $\Perm(n)$ は $|\Perm(n)| = n!$ の有限集合なので，集合
  \[
    \left\{ \frac{1}{n}\sum_{i=1}^{n} C_{i, \sigma(i)} \;\middle|\; \sigma \in \Perm(n) \right\}
  \]
  は高々 $n!$ 個の実数からなる空でない有限集合である．
  $\R$ の空でない有限部分集合は最小元を持つ（主張~\ref{clm:sem-finite-min}）
  ので，それを達成する $\sigma^* \in \Perm(n)$ が存在する．
\end{proof}

\begin{claim}{最適解が一意でない場合の存在}{sem-uniqueness}
  任意の $n \geq 2$ に対して，
  異なる $\sigma_1, \sigma_2 \in \Perm(n)$ がともに最小値を達成するような
  $\mathbf{C} \in \R_{\geq 0}^{n \times n}$ が存在する．
\end{claim}

\begin{proof}
  $C_{i,j} \defeq 1$（$i, j = 1, \ldots, n$）と定めると，任意の $\sigma \in \Perm(n)$ に対して
  $\frac{1}{n}\sum_{i=1}^{n} C_{i, \sigma(i)} = 1$．
  異なる $\sigma_1, \sigma_2 \in \Perm(n)$ を取れば，
  いずれも最小値 $1$ を達成する．
\end{proof}
